\section{Discussion}
\label{sec:discussion}

\textbf{What LAFC-Evict enables.} A reusable, candidate-level, finite-horizon
counterfactual supervision surface, materialized once from real traces and
usable across value-regression, best-candidate-selection, and pairwise
formulations (Sections~\ref{sec:background}--\ref{sec:tasks-baselines})
without every study re-deriving its own labels from raw traces
(Section~\ref{sec:related-work} situates this against prior learned-cache
work, most of which derives labels internally and non-reusably).

\textbf{What the evaluation reveals, beyond the surface itself.} Four
findings that only emerge from evaluating the surface against closed-loop
behavior, not from describing the surface alone: the target's
discriminativeness is workload-structured, not incidental
(Section~\ref{sec:characterization}); that structure is explained, not
just observed, by a single trace-level property, capacity-scale reuse
(Section~\ref{sec:mechanisms}); offline and closed-loop policy orderings
agree in most, not all, informative regimes, with the exceptions
identified rather than smoothed over (Section~\ref{sec:linkage}); and the
canonical LRU-continuation assumption is robust, with quantified limits,
at the capacities tested (Section~\ref{sec:continuation}).

\textbf{Why the negative/degenerate regimes are scientifically useful, not
an embarrassment.} wiki2018's complete degeneracy and alibaba-block/cap32's
near-zero offline signal are both, in this paper, explained results
(Section~\ref{sec:mechanisms}), not unexplained anomalies. A benchmark
whose degenerate cases are mechanistically understood is more trustworthy
in its non-degenerate claims than one that reports only favorable cases;
this is the paper's answer to the historical criticism that the target's
tie-heaviness was previously stated but not investigated.

\textbf{Offline surrogate evaluation vs.\ closed-loop performance.} This
paper's position, supported by Sections~\ref{sec:linkage}
and~\ref{sec:closed-loop}, is that the offline surrogate is a genuine,
partially-validated predictor of closed-loop policy separation in
identifiable regimes, and that this relationship is itself the
contribution -- not a claim that the surrogate can replace closed-loop
evaluation. The two exceptions found (alibaba-block/cap32,
metakv/cap128) are not edge cases invented for balance; they were
identified by the same symmetric procedure that found the 8 agreeing
cells, and metakv/cap128 in particular remains only partially understood
even after dedicated mechanistic analysis (Section~\ref{sec:mechanisms}).

\textbf{Continuation robustness as currently validated.} The robustness
finding in Section~\ref{sec:continuation} is deliberately reported with
its exact scope attached: sampled MRU and mean-random continuation at
capacities 32/128, and full-population MRU continuation at capacities
32/64/128/256. Within that scope, the LRU continuation used for label
generation is not driving the target ordering to a practically meaningful
degree: Set-C median Jaccard remains 1.0, mean cross-continuation regret
is below 0.001 misses in the population MRU census, and capacities 64/256
do not reveal hidden sensitivity. The result does not say that every
stateful deployed continuation policy would preserve the same ordering;
it says the tested alternatives did, with random tested only in the
validated sample.

\textbf{Methodological limitations that shape every result above} are
collected in Section~\ref{sec:limitations} rather than repeated here;
Section~\ref{sec:characterization}'s target-discriminativeness numbers,
Section~\ref{sec:linkage}'s small $n$, and Section~\ref{sec:continuation}'s
continuation-policy scope are the three limitations most load-bearing for
the claims in this discussion specifically.
